\documentclass[aspectratio=169]{beamer}
\usetheme{Madrid}
\usepackage[utf8]{inputenc}
\usepackage[T1]{fontenc}
\usepackage{tikz}
\usepackage{amssymb}
\usepackage{graphicx}
\setbeamertemplate{navigation symbols}{}
\setbeamercolor{title}{fg=white,bg=red!45!black}
\setbeamercolor{frametitle}{fg=white,bg=red!45!black}
\setbeamercolor{block title}{bg=red!45!black!20,fg=black}
\setbeamercolor{block body}{bg=red!45!black!8,fg=black}
\setbeamercolor{item}{fg=red!45!black}
\setbeamercolor{itemize item}{fg=red!45!black}
\setbeamertemplate{itemize item}[triangle]

\title{Bringing It Together}
\subtitle{Checkpoint — health, attack, win/lose, and turns, as one system}
\author{ByteForge --- Grade 8 --- App Builder}
\date{}

\begin{document}

\begin{frame}
    \titlepage
\end{frame}

\begin{frame}{Today's Plan}
    \begin{itemize}
        \item BOSS — Guided Mini-Battle
        \item[$\bigstar$] BOSS --- Guided Mini-Battle
    \end{itemize}
\end{frame}

\begin{frame}{A checkpoint combining cc1 through cc4, nothing new}
    \large
    \begin{itemize}
        \item A health variable, an attack that reduces it, a while loop that keeps fighting while health is above 0, and a win/lose message once it ends
        \item Genuinely just assembly of ideas already learned.
    \end{itemize}
\end{frame}


\begin{frame}{The outcome check happens AFTER the loop ends}
    \large
    \begin{itemize}
        \item \texttt{if health <= 0: print(...)} sits outside and after the \texttt{while} loop
        \item Checking the outcome only makes sense once the fighting has actually stopped.
    \end{itemize}
\end{frame}

\begin{frame}{Real-World Examples}
    \large
    \begin{itemize}
        \item This is genuinely the complete shape of the smallest possible real turn-based battle system
        \item The same four pieces (state, change, loop, outcome) scale up to a full commercial RPG's combat engine.
    \end{itemize}
\end{frame}


\begin{frame}{Let Them Discover}
    \Large
    Before opening the platform: on paper, list the four separate pieces a mini-battle needs (a starting health, a way to reduce it, a loop, an outcome check) — in what order do they have to happen?
\end{frame}

\begin{frame}{Your Turn}
    \Large
    Open the student map, find \textbf{Bringing It Together}, and work through it
    hands-on --- everything you just saw, they now do themselves.
\end{frame}

\end{document}
